% vim: ai sw=2

\documentclass[handout,xcolor=dvipsnames]{beamer}

\let\Oldemph\emph
\renewcommand{\emph}[1]{\textcolor{blue}{\Oldemph{#1}}}
\newcommand{\heading}[1]{\textbf{\textcolor{BurntOrange}{#1}}}

% Copyright 2004 by Till Tantau <tantau@users.sourceforge.net>.
%
% In principle, this file can be redistributed and/or modified under
% the terms of the GNU Public License, version 2.
%
% However, this file is supposed to be a template to be modified
% for your own needs. For this reason, if you use this file as a
% template and not specifically distribute it as part of a another
% package/program, I grant the extra permission to freely copy and
% modify this file as you see fit and even to delete this copyright
% notice. 


\mode<presentation>
{
%  \usetheme{Warsaw}
\usetheme{Madrid}

  \setbeamercovered{transparent}
}

\usepackage[utf8]{inputenc}
\usepackage{times}
\usepackage[T1]{fontenc}

\usepackage{graphicx}
\usepackage[final]{listings}
\usepackage{multirow}

\lstset{tabsize=2,numbers=none,showstringspaces=false,breaklines=true,
	breakatwhitespace=true,escapeinside={(*@}{@*)},
	basicstyle=\scriptsize,keywordstyle=\bfseries\color{OliveGreen},
	commentstyle=\color{blue}}

\title[Stanse] % (optional, use only with long paper titles)
{Stanse -- Taking a Firm Stanse on Bugs}

% \subtitle
% {Include Only If Paper Has a Subtitle}

\author[J. Obdržálek, J. Slabý, M. Trtík]{Jan~Obdržálek, \emph{Jiří~Slabý}, Marek~Trtík}

\institute[ITI] % (optional, but mostly needed)
{ITI, Faculty of Informatics, Masaryk University, Brno}

\date[03.03. 2010, Nuremberg]{Nuremberg, March 3 2010}

\subject{Theoretical Computer Science}
% This is only inserted into the PDF information catalog. Can be left
% out. 



% If you have a file called "university-logo-filename.xxx", where xxx
% is a graphic format that can be processed by latex or pdflatex,
% resp., then you can add a logo as follows:

% \pgfdeclareimage[height=0.5cm]{university-logo}{university-logo-filename}
% \logo{\pgfuseimage{university-logo}}



% Delete this, if you do not want the table of contents to pop up at
% the beginning of each subsection:
% \AtBeginSubsection[]
% {
%   \begin{frame}<beamer>
%     \frametitle{Outline}
%     \tableofcontents[currentsection,currentsubsection]
%   \end{frame}
% }


% If you wish to uncover everything in a step-wise fashion, uncomment
% the following command: 

%\beamerdefaultoverlayspecification{<+->}

\begin{document}

\begin{frame}
  \titlepage
\end{frame}

%%%%%%%%%%%%%%%%%%%%%%%

\begin{frame}
  \frametitle{Presentation outline}
  \begin{itemize}
  \item Background, the start-point.
  \item Features of the tool (configuration incl.).
  \item Found bugs (i.e. results).
  \item Work in progress and future work.
  \item Tool presentation \& discussion.
  \end{itemize}
\end{frame}

\begin{frame}
  \frametitle{Background}
  \heading{Formal verification, analysis, and testing of software
      systems}
  \begin{itemize}
  \item ITI/ANF Data Project (since 2006).
  \item Goal: automatically identify errors and bottlenecks in code.
  \item Aimed at real-life systems (large codebase etc.).
  \end{itemize}

  \pause
  \medskip
  \heading{Stanse}
  \begin{itemize}
  \item Developed at ITI, FI at the Masaryk University.
  \item First working version -- Summer 2007.
  \item Evaluated at ANF Data, found some bugs.
  \item Recently completely rewritten.
  \item Public release under GPLv2 in September 2009.
  \item \url{http://stanse.fi.muni.cz/}
  \end{itemize}

\end{frame}

\begin{frame}
  \frametitle{Stanse Features}
  \begin{itemize}
  \item Error-finding tool based on \emph{static analysis}.
  \item Target language is C (ANSI C99), but extensible to C\#/C++/Java.
  \item \emph{Full ANSI C99} support, including most GNU C extensions.
  \item Modular structure, easy extensibility, fast development.
  \item Easy to use \emph{graphical interface} and error path inspection.
  \item Common bug tracking system independent of checkers.
  \item Intraprocedural and interprocedural analyses.
  \item Batch execution + makefile support.
  \end{itemize}
\end{frame}

\begin{frame}
  \frametitle{Checkers}
  \heading{Automaton Checker}
  \begin{itemize}
  \item Uses \emph{state automata} to describe errorneous behavior.
  \item Finds locking, memory, reference count and other errors.
  \item Highly variable thanks to XML configuration.
  \begin{itemize}
    \item Basic userspace configurations provided.
    \item Kernel ones 
  \end{itemize}
  \end{itemize}
  \begin{center}
    \input{locking.pspdftex}
  \end{center}
\end{frame}
  
\begin{frame}
  \frametitle{Checkers contd.}
  \heading{Thread Checker}
  \begin{itemize}
  \item Finds threads and simultaneously executes them.
  \item Watches locks and unlocks (interprocedurally).
  \item Finds non-trivial locking problems across threads.
  \end{itemize}
  \bigskip
  \heading{Reachability Checker}
  \begin{itemize}
  \item Only 200 lines long.
  \item It serves as a tool demonstartion.
  \end{itemize}

\end{frame}

\begin{frame}[fragile]
  \frametitle{Configuration}
  \begin{verbatim}
  <bleble>
  \end{verbatim}
\end{frame}

\begin{frame}
\frametitle{Tools comparison}

\begin{tabular}{l|ccccl}
\multirow{2}{*}{Tool name} &
Open- &
Confi- &
Batch &
Inter- &
\multirow{2}{*}{Languages} \\
%Tool name        & Open & Confi-  & Batch   & Inter- & Languages
                  & source & gurable & support & proc.  & \\ \hline
\textsc{Coverity} & no   & yes     & yes     & yes    & GC, C++, C\#, Java \\
\textsc{Klocwork} & no   & yes     & yes     & yes    & GC, C++, C\#, Java \\
\textsc{Splint}   & yes  & no      & no      & no     & C \\
\textsc{Sparse}   & yes  & no      & yes     & no     & GC \\
\textsc{Clang}    & yes  & no      & no      & no     & GC$^*$, C++, Obj-C \\
\textsc{Uno}      & yes  & yes     & no      & yes    & C$^\dagger$ \\
\textsc{Stanse}   & yes  & yes     & yes     & yes    & GC \\
\textsc{FindBugs} & yes  & no      & yes     & no     & Java \\ \hline
\multicolumn{6}{l}{\footnotesize GC stands for GNU C.} \\
\multicolumn{6}{l}{\footnotesize $^*$\textsc{Clang} is still under heavy
development, its GNU C support is incomplete.} \\
\multicolumn{6}{l}{\footnotesize $^\dagger$C support in \textsc{Uno} does not
even implement full C99.} \\
\end{tabular}
\end{frame}

\begin{frame}
  \frametitle{What bugs have we found?}
  
  \heading{ANF Data}
  \begin{itemize}
  \item \texttt{malloc()} not tested for success
  \item unallocated memory freed
  \item \emph{memory leaks}
  \end{itemize}

  \medskip
  \heading{Linux kernel}\\
  Over 70 of the following kind of errors:
  \begin{itemize}
  \item unreachable code
  \item pointer usage violations
  \item locking discipline errors
    \begin{itemize}
    \item i.e. potential real-world deadlocks
    \end{itemize}
  \item function pairing
    \begin{itemize}
    \item e.g. get\_cpu/put\_cpu, interrupt disable/enable, etc.
    \end{itemize}
  \end{itemize}
\end{frame}

\begin{frame}[fragile]
  \frametitle{What bugs have we found?}

  One simple kernel example:
  \begin{lstlisting}[language=C]
static int set_lock_status(struct hotplug_slot *hotplug_slot, u8 status)
{
  struct slot *slot = hotplug_slot->private;
  ...
  mutex_lock(&slot->ctrl->crit_sect);

  /* has it been >1 sec since our last toggle? */
  if ((get_seconds() - slot->last_emi_toggle) < 1)
    return -EINVAL;
  ...
  mutex_unlock(&slot->ctrl->crit_sect);
  return 0;
}
  \end{lstlisting}

  \pause
  Full list at \url{http://stanse.fi.muni.cz/bugs.html}.

\end{frame}

\begin{frame}
  \frametitle{False Positives}
  \begin{itemize}
  \item Oscillating between 30 and 70\%.
  \item Stemming out of:
  \begin{itemize}
    \item Wrongly specified checkers.
    \begin{itemize}
      \item ble
    \end{itemize}
    \item All real errors evicted, only FP remains.
    \item Inappropriate (non-sophisticated) techniques.
  \end{itemize}
  \end{itemize}

\end{frame}

\begin{frame}[fragile]
  \frametitle{Fighting FP}
  \begin{itemize}
  \item Improving configuration.
  \begin{itemize}
    \item The quality is important and reflects direcly in FP ratio.
    \item Bugs in the configuration.
    \item Tune-up based on previous results.
  \end{itemize}
  
  \item FP detectors are java code finding certain patterns.
  \begin{itemize}
    \item For example tracking of condition \texttt{cond}:
  \end{itemize}
  \end{itemize}
  \begin{lstlisting}[language=C]
  if (cond) lock();
  if (cond) unlock();
  \end{lstlisting}

  \begin{itemize}
  \item And of course better analyses/checkers.
  \begin{itemize}
    \item E.g. those based on symbolic execution.
  \end{itemize}
  \end{itemize}

\end{frame}

\begin{frame}
  \frametitle{What can ITI offer?}
  \begin{itemize}
  \item BLE
  \end{itemize}

  \pause
  \vfill
  \begin{center}
    \emph{We want to hear about your problems! *}
  \end{center}

  \vfill
    {\tiny{* Applies only to errors in your (or somebody else's) code.}}
\end{frame}

\end{document}

